\section{Alternative Gate Sets and Expander Generators}

The all-transpositions model is attractive because it has rich Hurwitz
structure and very explicit local block conditionals.  The recent
polynomial-mixing work also suggests a possible drawback: the natural
$\{3,4,5\}$ heat-bath chain may have subtle slow modes inside weighted
support-component fibers or near sparse boundary layers.  It is therefore
reasonable to consider different generator sets.

A particularly clean theoretical alternative is to choose, for each $N$, a
bounded-size symmetric generating set $S_N^{\mathrm{exp}}\subseteq S_N$ or
$A_N$ whose Cayley graphs form an expander family.  Kassabov constructed such
bounded-size generating sets for symmetric and alternating groups; see
\href{https://arxiv.org/abs/math/0503204}{arXiv:math/0503204} and
\href{https://arxiv.org/abs/math/0505624}{arXiv:math/0505624}.  For the
unconditioned walk, this changes the baseline geometry: the Cayley graph has
bounded degree and a spectral gap bounded away from zero, so the ordinary
random walk mixes in $O(\log |S_N|)$ steps up to constants.

For identity-word sampling the target would become
\[
  \Omega_S(N,m)
  =
  \{(g_1,\ldots,g_m)\in S^m :
      g_1g_2\cdots g_m=\id\},
\]
with the uniform law on length-$m$ loops at the identity in
$\mathrm{Cay}(G_N,S_N^{\mathrm{exp}})$.  This is still a random-walk bridge
problem, but now on a known expander rather than on the complete
transposition Cayley graph equipped with a support-multigraph invariant.  A
natural research question is:
\[
  \text{Can bounded-degree expander Cayley graphs support efficient
  near-uniform loop samplers?}
\]
The expected proof tools would be different: bridge return-ratio estimates,
endpoint influence decay, and block dynamics on expander paths rather than
Hurwitz orbit classification and weighted support-component geometry.

There are two caveats.  First, expansion of the unconditioned Cayley walk does
not by itself prove rapid mixing of a local bridge sampler.  Endpoint
conditioning can introduce long-range dependence, and the exact fixed-length
bridge still requires either completion counts or a rapidly mixing
conditioned dynamics.  Second, Kassabov generators are global permutations of
$N$ points.  If $N=2^n$ is the state space of an $n$-wire reversible circuit,
such permutations are not automatically small local reversible gates.  They
are best viewed as a mathematically clean benchmark model unless an efficient
reversible-circuit implementation of the generators is supplied.

This suggests a useful split in future work:
\begin{itemize}
\item keep the all-transpositions model as the Hurwitz-rich model where exact
  local block samplers and support-component invariants are available;
\item study expander generator sets as a benchmark model for identity
  bridges on bounded-degree expanding Cayley graphs;
\item for cryptographic reversible circuits, compare both abstractions with
  implementable gate families from the random reversible-circuit literature,
  such as random local reversible gates or simple permutations.
\end{itemize}

If polynomial mixing for the $\{3,4,5\}$ transposition heat-bath chain is
false, an expander-generator bridge model could still provide a clean
theorem-level sampler target.  Conversely, if even expander-generator bridge
dynamics are hard to analyze, that would clarify that the main difficulty is
endpoint conditioning itself rather than the transposition support geometry.
